\subsection{BÀI TẬP TRẮC NGHIỆM}
\setcounter{ex}{0}
\ind{PHẦN I.} \inden{Câu trắc nghiệm nhiều phương án lựa chọn. Học sinh trả lời từ câu 1 đến câu 16. Mỗi câu hỏi học sinh chỉ chọn một phương án.}\\
\setcounter{ex}{0}
\Opensolutionfile{ans}[ans/1D2-B3-1]
\begin{ex}%[1D3Y4]
	Trong các dãy số hữu hạn sau, dãy số nào là cấp số nhân?
	\choice
	{$2; 4; 8; 16; 32; 63$}
	{\True $1; -2; 4; -8; 16; -32$}
	{$1; 3; 9; 27; 54; 162$} 
	{$4; 2; 1; \dfrac{1}{2}; \dfrac{1}{4}; \dfrac{1}{16}$}
	\loigiai{
		
		Vì $\dfrac{4}{2}= \dfrac{8}{4} = \dfrac{16}{8}= \dfrac{32}{16} \neq \dfrac{63}{32}$ nên dãy số ở phương án A không là cấp số nhân.
		
		Vì $\dfrac{3}{1}= \dfrac{9}{3} = \dfrac{27}{9} \ne \dfrac{54}{27}$ nên dãy số ở phương án C không là cấp số nhân.
		
		Vì $\dfrac{2}{4}= \dfrac{1}{2} = \dfrac{\frac{1}{2}}{1} = \dfrac{\frac{1}{4}}{\frac{1}{2}} \ne \dfrac{\frac{1}{16}}{\frac{1}{4}}$ nên dãy số ở phương án D không là cấp số nhân.
		
		Vì $\dfrac{-2}{1}= \dfrac{4}{-2} = \dfrac{-8}{4} = \dfrac{16}{-8} = \dfrac{-32}{16}=-2$ nên dãy số ở phương án B là cấp số nhân. 
		
	}
\end{ex}

\begin{ex}%[HK1, 2017 - 2018, Nguyễn Trung Ngạn, Hưng Yên]%[1D3B4]%
	Trong các đãy $(u_n)$ cho bởi số hạng tổng quát dưới đây, dãy nào là một cấp số nhân có công bội bằng $2$?
	\choice
	{$u_n = 2n + 3$}
	{\True $u_n = 2^n$}
	{$u_n = 2^n + 3$}
	{$u_n = n + 2$}
	\loigiai{
		$(u_n)$ là một cấp số nhân có công bội là $2$ $\Leftrightarrow$ $u_{n+1} = 2u_n$, $\forall n \in \mathbb{N}^*$. Kiểm tra các đáp án với vài số hạng đầu của dãy số thì ta thấy dãy số $(u_n)$ có số hạng tổng quát $u_n = 2^n$ là một cấp số nhân.
	}
\end{ex}

\begin{ex}%[HK1 lớp 11 - THPT Lộc Phát - Lâm Đồng]%[1D3B4]
	Tìm công bội của cấp số nhân $(u_n)$ biết số hạng tổng quát $u_n=3^{2n}$.
	\choice
	{\True $q=9$}
	{$q=2$}
	{$q=3$}
	{$q=6$}
\end{ex}


\begin{ex}
	Cấp số nhân $(u_n)$ có $u_1=3$, $q=2$. Tìm $u_2$.
	\choice
	{\True $u_2=6$}
	{$u_2=5$}
	{$u_2=-6$}
	{$u_2=1$}
	\loigiai{Cấp số nhân $(u_n)$ có $u_1=3$, $q=2$, có số hạng tổng quát $u_n=u_1 \cdot q^{n-1}$, $n \geq 2$. Vậy $u_2=u_1 \cdot q=6$.
	}
\end{ex}

\begin{ex}
	Cho cấp số nhân $(u_n)$ biết $u_1=3$ và công bội $q=-2$. Tìm $u_7$.
	\choice
	{\True $u_7=192$}
	{$u_7=-9$}
	{$u_7=-192$}
	{$u_7=384$}
	\loigiai
	{
		$u_7=u_1q^6=3(-2)^6=-192$.
	}
\end{ex}

\begin{ex}
	Cho cấp số nhân $\left({u_n}\right)$, biết $u_1=2,q=\dfrac{1}{3}$. Tìm $u_{10}$.
	\choice
	{$\dfrac{2}{3^8}$}
	{$\dfrac{2}{3^{10}}$}
	{$\dfrac{3}{2^9}$}
	{\True $\dfrac{2}{3^9}$}
	\loigiai{
		Cấp số nhân $\left({u_n}\right)$ có $u_1=2,q=\dfrac{1}{3} \Rightarrow u_{10}=u_1\cdot q^9=2\cdot \dfrac{1}{3^9}$.
	}
\end{ex}


\begin{ex}%[1D3Y4-5]
	Tính tổng 10 số hạng đầu tiên của cấp số nhân $(u_n)$, biết $u_1=-3$ và công bội $q=-2$.
	\choice
	{$S_{10}=-1023$}
	{$S_{10}=1025$}
	{$S_{10}=-1025$}
	{\True $S_{10}=1023$}
	\loigiai{
		Ta có: $S_{10}=\dfrac{u_1\left(1-q^{10}\right)}{1-q}=1023$.}
\end{ex}

\begin{ex}%[1D3Y4-5]%
	Một cấp số nhân $(u_n)$ có $u_1=-5$, $u_2=10$. Tổng của $15$ số hạng đầu của cấp số nhân đó là
	\choice
	{\True $2$}
	{$\dfrac{4}{3}$}
	{$\dfrac{2}{3}$}
	{$0$}
	\loigiai{
		Công bội của cấp số nhân đã cho là: $q=\dfrac{u_2}{u_1}=\dfrac{10}{-5}=-2$.\\
		Tổng của $15$ số hạng đầu của cấp số nhân đó là: $S_{15}=-5\cdot \dfrac{1-(-2)^{15}}{1-(-2)}=-54615$.
	}
\end{ex}

\begin{ex}%[1D3B4]
	Dãy số $\left(u_n\right)$ xác định bởi $\heva{&
		u_1 = 1\\&
		u_{n + 1} = \dfrac{1}{2}u_n}$ với $n \ge 1$. Tính tổng $S = u_1 + u_2 +\ldots + u_{10}$.
	\choice
	{$S = \dfrac{5}{2}$}
	{$S = \dfrac {1023} {2048}$}
	{$S = 2$}
	{\True $\dfrac {1023} {512}$}
	\loigiai{
		Ta có các số hạng của dãy số $\left(u_n\right)$ là  $1,\dfrac{1}{2},\dfrac{1}{4},\dfrac{1}{8},\dfrac {1}{16},\dfrac{1}{32},\ldots ,\dfrac{1} {2^n}$. Khi đó $\left(u_n\right)$ lập thành một cấp số nhân có $u_1 = 1$ và công bội $q = \dfrac{1}{2}$. \\
		Suy ra $S = u_1 + u_2 +\ldots + u_{10} = 1 + \dfrac{1}{2} + \dfrac {1}{4} +\ldots + \dfrac {1}{2^9}$$ = \dfrac {1\cdot \left(1 - \left( \dfrac{1}{2}\right)^{10}\right)} {1 - \dfrac{1}{2}} = \dfrac {1023} {512}$.}
\end{ex}

\begin{ex}%[1D3B4-5]
	Giá trị của tổng $S=1+3+3^2+\cdots+3^{2018}$ bằng
	\choice
	{\True $S=\dfrac{3^{2019}-1}{2}$}
	{$S=\dfrac{3^{2018}-1}{2}$}
	{$S=\dfrac{3^{2020}-1}{2}$}
	{$S=\dfrac{3^{2018}-1}{2}$}
	\loigiai{
		Ta thấy $S$ là tổng của $2019$ số hạng đầu tiên của cấp số nhân với số hạng đầu là $u_1=1$, công bội $q=3$.\\
		Áp dụng công thức tính tổng của cấp số nhân ta có $S=1\cdot\dfrac{1-3^{2019}}{1-3}=\dfrac{3^{2019}-1}{2}.$
	}  
\end{ex}   

\begin{ex}%[1D3B4-3]
	Cho cấp số nhân $(u_{n})$ có $u_{5}=15$ và $u_{8}=-1875$. Công bội của cấp số nhân là
	\choice
	{$q=3$}
	{$q=-3$}
	{\True $q=-5$}
	{$q=5$}
	\loigiai{
		Ta có $\heva{& u_5=u_1\cdot q^4=15\\& u_{8}=u_1\cdot q^7=-1875}\Leftrightarrow\heva{&u_1\cdot q^4=15\\&q^3=-125}\Leftrightarrow\heva{& u_1=\dfrac{3}{125}\\& q=-5.}$
	}
\end{ex}

\begin{ex}%[1D3B4-2]
	Cho cấp số nhân $(u_n)$ có $\heva{&u_{10}=8u_7 \\&u_1+u_4=144}$. Công bội $q$ của cấp số nhân là 
	\choice
	{$q=2$}
	{$q=-3$}
	{\True $q=3$}
	{$q=-2$}
	\loigiai
	{
		Ta có $u_{10}=8u_{7} \Leftrightarrow u_7\cdot q^3 =8u_7 \Leftrightarrow \hoac{&u_7=0\\&q=3}$.\\
		Vậy $q=3$.
	}
\end{ex}


\begin{ex}%[1D3Y4]
	Cho cấp số nhân với $ u_1 = 3, q = -2. $ Số $ 192 $ là số hạng thứ mấy của cấp số nhân?
	\choice
	{$ u_5 $}
	{$ u_6$}
	{\True $ u_7 $}
	{$ u_8 $}
	\loigiai{
		Ta có, $ u_n = u_1 . q^{n-1}  $ nên $ u_7 = 3 .(-2)^6 = 192. $
	}
\end{ex}

\begin{ex}%[1D3B4]
	Xác định $x$ để 3 số $2x-1$; $2x$; $2x+3$ theo thứ tự đó lập thành một cấp số nhân.
	\choice
	{$x=-\dfrac{3}{4}$}
	{$x=-\dfrac{4}{3}$}
	{\True $x=\dfrac{3}{4}$}
	{$x=\dfrac{4}{3}$}
\end{ex}

\begin{ex}%[1D3B4-3]
	Cho tam giác $ABC$ có ba góc $A,B,C$ theo thứ tự lập thành một cấp số nhân với công bội $q=2$. Tính số đo góc $A$.
	\choice
	{$\dfrac{\pi}{2}$}
	{$\dfrac{\pi}{7}$}
	{\True $\dfrac{2\pi}{7}$}
	{$\dfrac{4\pi}{7}$}
	\loigiai{
		Theo giả thiết ta có $A,B,C$ lập thành cấp số nhân với $q=2$ nên $B=2A;C=4A$ \\
		Mà $ABC$ là một tam giác nên $A+2A+4A=\pi \Leftrightarrow A=\dfrac{\pi}{7}$.}
\end{ex}


\begin{ex}%[1D3G4-7]
	Một thợ thủ công muốn vẽ trang trí một hình vuông kích thước $ 4\, \mathrm{m}\times 4\, \mathrm{m} $ bằng cách vẽ một hình vuông mới với các đỉnh là trung điểm các cạnh của hình vuông ban đầu, và tô kín màu lên hai tam giác đối diện (như hình vẽ). Quá trình vẽ và tô theo quy luật đó được lặp lại $ 5 $ lần. Tính số tiền nước sơn để người thợ đó hoàn thành trang trí hình vuông trên? Biết tiền nước sơn $ 1\, \mathrm{m}^2 $ là $ 60\,000 $ đồng.
	\begin{center}
		{\begin{tikzpicture}[scale=0.8,>=stealth, font=\footnotesize, line join=round, line cap=round]
				\def\xmin{-2.5}\def\xmax{2.5}\def\ymin{-2.5}\def\ymax{2.5} 
				%\draw[color=gray!50,dashed] (\xmin,\ymin) grid (\xmax,\ymax); 
				%\draw[->] (\xmin,0)--(\xmax,0) node [below]{$x$}; 
				%\draw[->] (0,\ymin)--(0,\ymax) node [left]{$y$}; 
				\clip (\xmin,\ymin) rectangle (\xmax,\ymax); 
				%%%%
				\def\r{4}
				\coordinate (O) at (0,0); 
				\foreach\n in {2}
				{
					\coordinate(A) at ($(O)+({(2^(-\n/2))*\r*cos(-\n*45)},{(2^(-\n/2))*\r*sin(-\n*45)})$); 
					\coordinate(B) at ($(O)+({(2^(-\n/2))*\r*cos(-\n*45-90)},{(2^(-\n/2))*\r*sin(-\n*45-90)})$); 
					\coordinate(C) at ($(O)+({(2^(-\n/2))*\r*cos(-\n*45-180)},{(2^(-\n/2))*\r*sin((-\n*45-180)})$); 
					\coordinate(D) at ($(O)+({(2^(-\n/2))*\r*cos(-\n*45-270)},{(2^(-\n/2))*\r*sin(-\n*45-270)})$); 
					\fill[blue!30] (C)--(D)--($(O)+({(2^(-(\n-1)/2))*\r*cos(-(\n-1)*45-270)},{(2^(-(\n-1)/2))*\r*sin(-(\n-1)*45-270)})$)--cycle; 
					\fill[pink] (A)--(B)--($(O)+({(2^(-(\n-1)/2))*\r*cos(-(\n-1)*45-90)},{(2^(-(\n-1)/2))*\r*sin(-(\n-1)*45-90)})$)--cycle; 
					\draw (A)--(B)--(C)--(D)--cycle; 
				}
				\def\n{1}
				\coordinate(A) at ($(O)+({(2^(-\n/2))*\r*cos(-\n*45)},{(2^(-\n/2))*\r*sin(-\n*45)})$); 
				\coordinate(B) at ($(O)+({(2^(-\n/2))*\r*cos(-\n*45-90)},{(2^(-\n/2))*\r*sin(-\n*45-90)})$); 
				\coordinate(C) at ($(O)+({(2^(-\n/2))*\r*cos(-\n*45-180)},{(2^(-\n/2))*\r*sin((-\n*45-180)})$); 
				\coordinate(D) at ($(O)+({(2^(-\n/2))*\r*cos(-\n*45-270)},{(2^(-\n/2))*\r*sin(-\n*45-270)})$); 
				\draw (A)--(B)--(C)--(D)--cycle; 
		\end{tikzpicture}}
		{\begin{tikzpicture}[scale=0.8,>=stealth, font=\footnotesize, line join=round, line cap=round]
				\def\xmin{-2.5}\def\xmax{2.5}\def\ymin{-2.5}\def\ymax{2.5} 
				%\draw[color=gray!50,dashed] (\xmin,\ymin) grid (\xmax,\ymax); 
				%\draw[->] (\xmin,0)--(\xmax,0) node [below]{$x$}; 
				%\draw[->] (0,\ymin)--(0,\ymax) node [left]{$y$}; 
				\clip (\xmin,\ymin) rectangle (\xmax,\ymax); 
				%%%%
				\def\r{4}
				\coordinate (O) at (0,0); 
				\foreach\n in {2,3}
				{
					\coordinate(A) at ($(O)+({(2^(-\n/2))*\r*cos(-\n*45)},{(2^(-\n/2))*\r*sin(-\n*45)})$); 
					\coordinate(B) at ($(O)+({(2^(-\n/2))*\r*cos(-\n*45-90)},{(2^(-\n/2))*\r*sin(-\n*45-90)})$); 
					\coordinate(C) at ($(O)+({(2^(-\n/2))*\r*cos(-\n*45-180)},{(2^(-\n/2))*\r*sin((-\n*45-180)})$); 
					\coordinate(D) at ($(O)+({(2^(-\n/2))*\r*cos(-\n*45-270)},{(2^(-\n/2))*\r*sin(-\n*45-270)})$); 
					\fill[blue!30] (C)--(D)--($(O)+({(2^(-(\n-1)/2))*\r*cos(-(\n-1)*45-270)},{(2^(-(\n-1)/2))*\r*sin(-(\n-1)*45-270)})$)--cycle; 
					\fill[pink] (A)--(B)--($(O)+({(2^(-(\n-1)/2))*\r*cos(-(\n-1)*45-90)},{(2^(-(\n-1)/2))*\r*sin(-(\n-1)*45-90)})$)--cycle; 
					\draw (A)--(B)--(C)--(D)--cycle; 
				}
				\def\n{1}
				\coordinate(A) at ($(O)+({(2^(-\n/2))*\r*cos(-\n*45)},{(2^(-\n/2))*\r*sin(-\n*45)})$); 
				\coordinate(B) at ($(O)+({(2^(-\n/2))*\r*cos(-\n*45-90)},{(2^(-\n/2))*\r*sin(-\n*45-90)})$); 
				\coordinate(C) at ($(O)+({(2^(-\n/2))*\r*cos(-\n*45-180)},{(2^(-\n/2))*\r*sin((-\n*45-180)})$); 
				\coordinate(D) at ($(O)+({(2^(-\n/2))*\r*cos(-\n*45-270)},{(2^(-\n/2))*\r*sin(-\n*45-270)})$); 
				\draw (A)--(B)--(C)--(D)--cycle; 
		\end{tikzpicture}}
		{\begin{tikzpicture}[scale=0.8,>=stealth, font=\footnotesize, line join=round, line cap=round]
				\def\xmin{-2.5}\def\xmax{2.5}\def\ymin{-2.5}\def\ymax{2.5} 
				%\draw[color=gray!50,dashed] (\xmin,\ymin) grid (\xmax,\ymax); 
				%\draw[->] (\xmin,0)--(\xmax,0) node [below]{$x$}; 
				%\draw[->] (0,\ymin)--(0,\ymax) node [left]{$y$}; 
				\clip (\xmin,\ymin) rectangle (\xmax,\ymax); 
				%%%%
				\def\r{4}
				\coordinate (O) at (0,0); 
				\foreach\n in {2,3,4}
				{
					\coordinate(A) at ($(O)+({(2^(-\n/2))*\r*cos(-\n*45)},{(2^(-\n/2))*\r*sin(-\n*45)})$); 
					\coordinate(B) at ($(O)+({(2^(-\n/2))*\r*cos(-\n*45-90)},{(2^(-\n/2))*\r*sin(-\n*45-90)})$); 
					\coordinate(C) at ($(O)+({(2^(-\n/2))*\r*cos(-\n*45-180)},{(2^(-\n/2))*\r*sin((-\n*45-180)})$); 
					\coordinate(D) at ($(O)+({(2^(-\n/2))*\r*cos(-\n*45-270)},{(2^(-\n/2))*\r*sin(-\n*45-270)})$); 
					\fill[blue!30] (C)--(D)--($(O)+({(2^(-(\n-1)/2))*\r*cos(-(\n-1)*45-270)},{(2^(-(\n-1)/2))*\r*sin(-(\n-1)*45-270)})$)--cycle; 
					\fill[pink] (A)--(B)--($(O)+({(2^(-(\n-1)/2))*\r*cos(-(\n-1)*45-90)},{(2^(-(\n-1)/2))*\r*sin(-(\n-1)*45-90)})$)--cycle; 
					\draw (A)--(B)--(C)--(D)--cycle; 
				}
				\def\n{1}
				\coordinate(A) at ($(O)+({(2^(-\n/2))*\r*cos(-\n*45)},{(2^(-\n/2))*\r*sin(-\n*45)})$); 
				\coordinate(B) at ($(O)+({(2^(-\n/2))*\r*cos(-\n*45-90)},{(2^(-\n/2))*\r*sin(-\n*45-90)})$); 
				\coordinate(C) at ($(O)+({(2^(-\n/2))*\r*cos(-\n*45-180)},{(2^(-\n/2))*\r*sin((-\n*45-180)})$); 
				\coordinate(D) at ($(O)+({(2^(-\n/2))*\r*cos(-\n*45-270)},{(2^(-\n/2))*\r*sin(-\n*45-270)})$); 
				\draw (A)--(B)--(C)--(D)--cycle; 
		\end{tikzpicture}}
		{\begin{tikzpicture}[scale=0.8,>=stealth, font=\footnotesize, line join=round, line cap=round]
				\def\xmin{-2.5}\def\xmax{2.5}\def\ymin{-2.5}\def\ymax{2.5} 
				%\draw[color=gray!50,dashed] (\xmin,\ymin) grid (\xmax,\ymax); 
				%\draw[->] (\xmin,0)--(\xmax,0) node [below]{$x$}; 
				%\draw[->] (0,\ymin)--(0,\ymax) node [left]{$y$}; 
				\clip (\xmin,\ymin) rectangle (\xmax,\ymax); 
				%%%%
				\def\r{4}
				\coordinate (O) at (0,0); 
				\foreach\n in {2,3,4,5}
				{
					\coordinate(A) at ($(O)+({(2^(-\n/2))*\r*cos(-\n*45)},{(2^(-\n/2))*\r*sin(-\n*45)})$); 
					\coordinate(B) at ($(O)+({(2^(-\n/2))*\r*cos(-\n*45-90)},{(2^(-\n/2))*\r*sin(-\n*45-90)})$); 
					\coordinate(C) at ($(O)+({(2^(-\n/2))*\r*cos(-\n*45-180)},{(2^(-\n/2))*\r*sin((-\n*45-180)})$); 
					\coordinate(D) at ($(O)+({(2^(-\n/2))*\r*cos(-\n*45-270)},{(2^(-\n/2))*\r*sin(-\n*45-270)})$); 
					\fill[blue!30] (C)--(D)--($(O)+({(2^(-(\n-1)/2))*\r*cos(-(\n-1)*45-270)},{(2^(-(\n-1)/2))*\r*sin(-(\n-1)*45-270)})$)--cycle; 
					\fill[pink] (A)--(B)--($(O)+({(2^(-(\n-1)/2))*\r*cos(-(\n-1)*45-90)},{(2^(-(\n-1)/2))*\r*sin(-(\n-1)*45-90)})$)--cycle; 
					\draw (A)--(B)--(C)--(D)--cycle; 
				}
				\def\n{1}
				\coordinate(A) at ($(O)+({(2^(-\n/2))*\r*cos(-\n*45)},{(2^(-\n/2))*\r*sin(-\n*45)})$); 
				\coordinate(B) at ($(O)+({(2^(-\n/2))*\r*cos(-\n*45-90)},{(2^(-\n/2))*\r*sin(-\n*45-90)})$); 
				\coordinate(C) at ($(O)+({(2^(-\n/2))*\r*cos(-\n*45-180)},{(2^(-\n/2))*\r*sin((-\n*45-180)})$); 
				\coordinate(D) at ($(O)+({(2^(-\n/2))*\r*cos(-\n*45-270)},{(2^(-\n/2))*\r*sin(-\n*45-270)})$); 
				\draw (A)--(B)--(C)--(D)--cycle; 
		\end{tikzpicture}}
	\end{center}
	\choice
	{$575\,000 $ đồng}
	{$387\,500 $ đồng}
	{\True $465\,000 $ đồng}
	{$232\,500 $ đồng}
	\loigiai{
		Gọi $ a $ là độ dài cạnh hình vuông. Độ dài cạnh hình vuông theo lần vẽ thứ $ n $ theo quy luật trên là $ 2^{-\frac{n}{2}}a $. Tổng diện tích của hai tam giác được tô màu sau lần vẽ thứ $ n $ theo quy luật trên là $ S_n=\dfrac{1}{2}\left(2^{-\frac{n}{2}}\right)^2a^2=2^{-n-1}a^2. $\\
		Gọi $ u_n=S_1+S_2+\cdots+S_n=2^{-2}a^2+2^{-3}a^2+\cdots+2^{-n-1}a^2=2^{-1}\left(1-\dfrac{1}{2^n}\right)a^2 $.\\
		Theo đề bài ta có $ n=5 $ và $ a=4 $ nên $ u_5=\dfrac{31}{4}\ m^2 $.\\
		Vậy số tiền nước sơn là $\dfrac{31}{4}\cdot 60000=465000 $ đồng.
	}
\end{ex}

\Closesolutionfile{ans}

\ind{PHẦN II.} \inden{Câu trắc nghiệm đúng sai. Học sinh trả lời từ câu 17 đến câu 20. Trong mỗi ý a), b), c), d) ở mỗi câu, học sinh chọn đúng hoặc sai.}\\
\Opensolutionfile{ans}[ans/1D2-B3-2]

\begin{ex}%[1C2Y3-3]
	Cho cấp số nhân $(u_n)$ với số hạng đầu $u_1 = -5$, công bội $q = 2$. Xét tính đúng sai của các khẳng định sau:
	\choiceTF
	{$u_9=1280$}
	{$u_n =-5\cdot 2^{n}$}
	{$160$ là một số hạng của cấp số nhân trên}
	{\True Tổng của 10 số hạng đầu bằng $-5115$}
	\loigiai{
		\begin{itemchoice}
			\itemch $u_9 = -5 \cdot 2^8 = -1280$.
			\itemch $u_n = u_1 \cdot q^{n-1} = (-5) \cdot 2^{n-1}$
			\itemch Do $u_n = u_1 \cdot q^{n-1} = (-5) \cdot 2^{n-1} < 0$ với mọi $n \in \mathbb{N}^*$ nên số $160$ không là số hạng nào của $(u_n)$.
			\itemch $S_{10}=-5 \cdot \dfrac{1-2^{10}}{1-2}=-5115$.
		\end{itemchoice}
	}
\end{ex}

\begin{ex}%[1D2B3-5]
	Cho một cấp số nhân $(u_n)$ có các số hạng đều không âm và thoả mãn $u_2=6$ và $u_4=24$. Xét tính đúng sai của các khẳng định sau:
	\choiceTF
	{Công bội của cấp số nhân là $q= \pm 2$}
	{\True Số hạng thứ ba của cấp số nhân bằng $12$}
	{\True $S_{12}=12285$}
	{\True $3072$ là một số hạng của cấp số nhân}
	\loigiai{
		Gọi công bội của cấp số nhân là $q$.\\
		Ta có $u_4=u_2\cdot q^2$ hay $24=6\cdot q^2$, suy ra $q=-2$ hoặc $q=2$. Do cấp số nhân có các số hạng không âm nên $q=2$.\\
		Khi đó, cấp số nhân $(u_n)$ có $u_1=\dfrac{u_2}{q}=\dfrac{6}{2}=3$.
		\begin{itemchoice}
			\itemch $q=2$
			\itemch $u_3=u_1.q^2=3 \cdot 4 =12$.
			\itemch $S_{12}=\dfrac{u_1\left(1-q^{12}\right)}{1-q}=\dfrac{3\left(1-2^{12}\right)}{1-2}=12285$.
			\itemch Số hạng tổng quát $u_n=3.2^{n-1}=3072 \Rightarrow 2^{n-1}=1024 \Rightarrow n=11$. Vậy $3072$ là số hạng thứ 11 của cấp số nhân
		\end{itemchoice}
	}
\end{ex}

\begin{ex}%[1D2V3-3]
	Cho dãy số $(u_n)$, biết $u_1=8$, $u_{n+1}=4u_n-9$ với $n \in \mathbb{N^*}$. Đặt $v_n=u_n-3$ với $n \in \mathbb{N^*}$.
	\choiceTF
	{$v_1=5$}
	{Dãy số $(v_n)$ là một cấp số nhân có công bội $q=-3$}
	{Công thức của số hạng tổng quát $v_n$ là $v_n=5 \cdot (-3)^{n-1}$}
	{Công thức của số hạng tổng quát $u_n$ là $u_n=3+5 \cdot (-3)^{n-1}$}
	\loigiai{
		Xét
		\begin{eqnarray*}
			&&u_{n+1}=4u_n-9 \Leftrightarrow u_n=4 \cdot u_{n-1}-9\\
			&\Rightarrow&v_n=u_n-3=4u_{n-1}-12\\
			&\Rightarrow&v_n=4(v_{n-1}+3)-12\\
			&\Rightarrow&4v_{n-1}+12-12=4v_{n-1}.
		\end{eqnarray*}
		Ta có $u_1=4u_0-9 \Leftrightarrow 8=4u_0-9 \Rightarrow u_0=\dfrac{17}{4}$. Mà $v_0=u_0-3 \Rightarrow v_0=\dfrac{17}{4}-3=\dfrac{5}{4}$.
		\begin{itemchoice}
			\itemch Đúng. Vì $v_1=4v_0=4 \cdot \dfrac{5}{4}=5$.
			\itemch Sai. Vì $v_n=4v_{n-1}$ là cấp số nhân với công bội bằng $4$.
			\itemch Sai. Vì $v_n=u_1 \cdot q^{n-1}=5 \cdot 4^{n-1}$.
			\itemch Sai. Vì $u_n=v_n+3=5 \cdot 4^{n-1}+3$.
		\end{itemchoice}
	}
\end{ex}

\begin{ex}%[1D2V3-4]
	Cho dãy số $(u_n)$ có tổng $n$ số hạng đầu được tính bởi công thức $S_n=\dfrac{1-3^n}{2\cdot 3^{n-2}}$ với $n \in \mathbb{N^*}$. Xét tính đúng sai của các khẳng định sau:
	\choiceTF
	{\True Số hạng thứ nhất của dãy là $u_1=-3$}
	{Số hạng thứ hai của dãy là $u_2=-4$}
	{Số hạng tổng quát của dãy số là $u_n=\dfrac{1}{3^{n-2}}$}
	{Dãy số $(u_n)$ là một cấp số nhân có công bội là $-\dfrac{1}{3}$}
	\loigiai{
		\begin{itemchoice}
			\itemch Đúng. Vì $u_1=S_1=\dfrac{1-3^1}{2\cdot 3^{1-2}}=-3$.
			\itemch Sai. Vì $S_2=u_1+u_2=-4$ nên $u_2=S_2-S_1=-1$.
			\itemch Sai. Vì ta có $S_n = u_1 + u_2 + \cdots + u_n=S_{n-1}+u_n$ nên $u_n = S_n - S_{n-1}$.\\
			Đầu tiên, Ta có $S_{n-1} = \dfrac{1 - 3^{n-1}}{2 \cdot 3^{n-3}}= \dfrac{3 \cdot (1 - 3^{n-1})}{3\cdot 2 \cdot 3^{n-3}} = \dfrac{3 - 3^n}{2 \cdot 3^{n-2}}$ và $S_n = \dfrac{1 - 3^n}{2 \cdot 3^{n-2}}$.\\
			Suy ra $\begin{aligned}[t]
				u_n &= S_n - S_{n-1}\\
				&= \dfrac{1 - 3^n}{2 \cdot 3^{n-2}} - \dfrac{3 - 3^n}{2 \cdot 3^{n-2}} \\
				&= \dfrac{1 - 3 - (3 - 3^n)}{2 \cdot 3^{n-2}} \\
				&= \dfrac{-2}{2 \cdot 3^{n-2}} \\
				&= -\dfrac{1}{3^{n-2}}.
			\end{aligned}$\\
			Điều này cho thấy $u_n = -\dfrac{1}{3^{n-2}}$.
			\itemch Sai. Vì ta có $u_n=u_{n-1}\cdot \dfrac{1}{3}$ với mọi $n\geq2$. Vậy nên dãy số $(u_n)$ là một cấp số nhân có số hạng đầu là $u_1=-3$ và cộng bội $q=\dfrac{1}{3}$.
		\end{itemchoice}
	}
\end{ex}

\Closesolutionfile{ans}

\ind{PHẦN III.} \inden{Câu trắc nghiệm trả lời ngắn. Học sinh trả lời từ câu 21 đến câu 25 vào ô kết quả.}\\
\Opensolutionfile{ans}[ans/1D2-B3-3]

\begin{ex}
	Cho cấp số nhân $(u_n)$ có số hạng đầu $u_1 = 2$ và công bội $q = -2$. Giá trị $u_5$ bằng bao nhiêu?
	\shortans{$32$}
	\loigiai{
		Ta có $u_5 = u_1\cdot q^4 = 2 \cdot (-2)^4 = 32$.	
	}
\end{ex}

\begin{ex}%[Nguyễn Đắc Giáp]%[1D3Y4-5]%
	Một cấp số nhân có số hạng đầu $u_1=3$ và công bội $q=2 $. Tổng $8$ số hạng đầu tiên của cấp số nhân bằng bao nhiêu?\\
	\shortans{$765$}
	\loigiai{
		Ta có
		$$S_8=\dfrac{u_1\left(1-q^8\right)}{1-q}=\dfrac{3 \left(1-2^8\right)}{1-2}=765.$$
	}
\end{ex}

\begin{ex}%Câu 42%[Trương Hữu Đăng]%[1C2K3-3]
	Ba số phân biệt tạo thành một cấp số nhân có tổng bằng $78$; đồng thời chúng là số hạng thứ nhất, thứ ba và thứ chín của một cấp số cộng. Tính tổng của ba số đó.\\
	\shortans{$78$}
	\loigiai{
		Giả sử công bội của cấp số nhân là $q$, công sai của cấp số cộng là $d$, khi đó gọi ba số cần tìm là $a$, $aq$, $aq^2$ (với $a$, $p \neq 0$).\\
		Theo bài ra ta có $\heva{&a + aq + aq^2 = 78 \;\; (*)\\&aq = a + 2d\\&aq^2 = a + 8d.}$\\
		Từ $aq = a + 2d$, suy ra $aq - a = 2d \Leftrightarrow a(q - 1) = 2d$. (1)\\
		Từ $aq^2 = a + 8d$, suy ra $$aq^2 - a = 8d \Leftrightarrow a\left(q^2 - 1\right) = 8d \Leftrightarrow a(q - 1)(q + 1) = 8d. (2)$$ 
		Với $q = 1$ thì $a = aq = aq^2$, mà ba số cần tìm là phân biệt nên $q = 1$ không thỏa mãn.\\
		Do vậy, $q \neq 1 \Rightarrow q-1 \neq 0$, do đó $a(q - 1) \neq 0$.\\
		Chia vế theo vế của (2) cho (1) ta được $q + 1 = 4 \Leftrightarrow q = 3$.\\
		Thay $q = 3$ vào (*) ta được
		Z$$a + 3a + 9a = 78 \Leftrightarrow 13a = 78 \Leftrightarrow a = 6.$$
		Suy ra ba số cần tìm là $6$; $6 \cdot 3 = 18$; $18 \cdot 3 = 54$.\\
		Vậy ba số cần tìm là $6$; $18$; $54$. Tổng ba số bằng $78$.	
	}
\end{ex}

\begin{ex}%[1D2V3-8]
	Một quả bóng được thả thẳng đứng từ độ cao $10$m rơi xuống đất và nảy lên. Giả sử sau mối một lần rơi xuống, nó nảy lên được một độ cao bằng $75\%$ độ cao vừa rơi xuống. Tính tổng quãng đường quả bóng di chuyển được kể từ lúc thả xuống đến khi quả bóng chạm đất lần thứ $10$ (làm tròn kết quả đến hàng phần mười của mét).\\
	\shortans{$65{,}5$}
	\loigiai{
		Gọi $u_n$(m) là độ cao mà quả bóng đạt được sau khi nảy lên ở lần thứ $n$. \\
		Ta có $u_1=10\cdot 0{,}75=7{,}5$. Ta có, dãy $(u_n)$ lập thành cấp số nhân có $u_1=7{,}5$ và công bội $q=0{,}75$. Kể từ lúc thả xuống đến khi quả bóng chạm đất lần thứ $10$, quả bóng đã được nảy lên $9$ lần rồi lại rơi xuống.\\
		Do quãng đường quả bóng nảy lên và rơi xuống bằng nhau nên tổng quãng đường quả bóng di chuyển được kể từ lúc thả xuống đến khi quả bóng chạm đất lần thứ $10$ là
		$$
		S=10+2\left(u_1+u_2+\cdots+u_9\right)=10+2 \cdot 7{,}5 \cdot \dfrac{1-(0{,}75)^9}{1-0{,}75} \approx 65{,}5\text{ (m)}.
		$$
	}
\end{ex}

\begin{ex}%Câu 45%[Trương Hữu Đăng]%[1C2K3-7]
	Anh Dũng kí hợp đồng lao động trong 10 năm với phương án trả lương như sau: Năm thứ nhất, tiền lương của anh Dũng là $120$ triệu đồng. Kể từ năm thứ hai trở đi, mỗi năm tiền lương của anh Dũng được tăng lên $10\%$. Tính tổng số tiền lương anh Dũng lĩnh được trong $10$ năm đầu đi làm (làm tròn kết quả đến hàng đơn vị theo đơn vị triệu đồng).\\
	\shortans{$1912$}
	\loigiai{
		Ta có tiền lương năm thứ nhất của anh Dũng là $120$ triệu đồng.\\
		Tiền lương năm thứ hai của anh Dũng là
		$120 + 120 \cdot 10 \% = 120(1 + 0{,}1) = 120 \cdot 1{,}1$ (triệu đồng).\\
		Tiền lương năm thứ ba của anh Dũng là
		$$120 \cdot 1{,}1 + 120 \cdot 1{,}1 \cdot 10 \% = 120 \cdot 1{,}1(1 + 0{,}1) = 120 \cdot 1{,}1^2 \;\; \text{(triệu đồng).}$$
		Cứ tiếp tục như vậy, ta được tiền lương năm thứ 10 của anh Dũng là $120 \cdot 1{,}19$ (triệu đồng).\\
		Do vậy, tiền lương mỗi năm của anh Dũng nhận được trong $10$ năm lập thành một cấp số nhân với số hạng đầu $u_1=120$ và công bội $q = 1{,}1$.\\
		Khi đó tổng số tiền lương anh Dũng lĩnh được trong 10 năm đầu đi làm là
		$$S_{10} = \dfrac{u_1 \left(1 - q^{10}\right)}{1-q} = \dfrac{120\left[1 - (1{,}1)^{10}\right]}{1 - 1{,}1} \approx 1912 \;\; \text{(triệu đồng).}
		$$
		Vậy tổng số tiền lương anh Dũng lĩnh được trong $10$ năm đầu đi làm là $1912$ triệu đồng.	
	}
\end{ex}


\Closesolutionfile{ans}



